It is important to establish the context in which we will want to study photons. One of the novel ways researchers use the LHC in the last few years is as a Large Photon Collider. This is done by utilizing both pp, pA and A + A ultra peripheral collisions as sources of photons due to high acceleration of charged particles, combined with lack of hadronic interactions due to geometry constraints. As a result, one can measure extremely high energy photon-photon and photon-matter collisions, making it a very exciting probe for a variety of questions both in nuclear and high energy physics \cite{strikman_probing_2006,iancu_probing_2023,mantysaari_probing_2023}. Therefore, the following section will describe photon interactions under the lens of ``photon colliders''. 